\chapter{Functions of Several Variables: Surfaces and Level Curves}

A function of one variable is understood by looking at its graph. For a
function of two variables there are two pictures, and both are needed. The
\textbf{graph} is a surface in space: mountains, valleys, passes. The
\textbf{contour map} is what a cartographer draws instead --- the level
curves of that landscape, projected flat onto the plane. Economics almost
never draws the surface. It draws the contour map and gives the curves
names: isoquants, indifference curves, iso-profit lines.

The reason to spend a chapter on pictures before computing anything is that
most of multivariable calculus is easier to see than to calculate. Whether a
critical point is a maximum, a minimum or a saddle is visible at a glance in
a contour map. So is whether a function is convex, whether its upper contour
sets are convex, and where it is steep. And the geometry is not decoration:
the Lagrange method is the statement that two families of level curves touch
without crossing, and Kuhn--Tucker is a statement about which boundary of a
shaded region binds.

This chapter therefore does two things. It sets up domains, graphs, level
curves and the sub- and superlevel sets, with the vocabulary in both the
German and the English forms; and it drills reading properties off a picture
in both directions --- picture to property, and formula to picture.

\section{Core Rules}

\begin{rulebox}{Functions of two variables: domain, graph, range}
A function $f\colon D\to\R$ with $D\subseteq\R^{2}$ assigns a number
$f(x,y)$ to each point of a \textbf{region} $D$ in the plane.
\begin{itemize}[leftmargin=1.6em,itemsep=0.3ex]
  \item The \textbf{domain} $D$ is now a subset of the plane, not of the
        line, and the same three obstructions decide it: no zero
        denominator, no negative radicand under an even root, no
        non-positive argument in a logarithm. So
        $f(x,y)=\sqrt{9-x^{2}-y^{2}}$ has $D=\{(x,y):x^{2}+y^{2}\le9\}$, a
        closed disc, and $f(x,y)=\ln(x+y)$ has the open half-plane
        $x+y>0$.
  \item The \textbf{graph} is the surface
        $\{(x,y,f(x,y)) : (x,y)\in D\}$ in $\R^{3}$: above each point of
        the domain, a single height.
  \item The \textbf{range} is the set of heights attained, exactly as
        before.
\end{itemize}
Economic domains are usually the non-negative quadrant $K,L\ge0$, and often
its interior.
\end{rulebox}

\begin{rulebox}{Level curves and the contour map}
For $c\in\R$ the \textbf{level set} \de{(Niveaumenge, Niveaulinie)} is
\[
  N_{c}=\{(x,y)\in D : f(x,y)=c\} ,
\]
the set of points at height $c$. Drawing $N_{c}$ for several equally spaced
values of $c$ gives the \textbf{contour map}. To compute one, set
$f(x,y)=c$ and solve for $y$ --- the level set is a curve in the plane, and
often a familiar one.

Reading a contour map:
\begin{itemize}[leftmargin=1.6em,itemsep=0.3ex]
  \item \textbf{Contours cannot cross}, since a point has one height ---
        with one exception: at a \textbf{saddle} the level curve through
        the saddle point crosses itself.
  \item \textbf{Spacing is steepness.} Contours crowded together mean the
        surface rises fast; widely spaced contours mean gentle ground.
        Steepest ascent runs perpendicular to the contours.
  \item \textbf{Closed nested loops} shrinking to a point mark a
        \textbf{peak} or a \textbf{basin}; which one is decided by whether
        the labels increase or decrease inwards.
  \item A contour that \textbf{pinches and crosses} marks a saddle: two
        directions of ascent and two of descent alternate around it.
\end{itemize}
\end{rulebox}

\begin{rulebox}{Sub- and superlevel sets}
Two \emph{solid} sets accompany each level curve:
\[
  W(c)=\{(x,y): f(x,y)\le c\},
  \qquad
  B(c)=\{(x,y): f(x,y)\ge c\} .
\]
$W(c)$ is the \textbf{sublevel set} \de{(Subniveaumenge, untere
Konturmenge)}, $B(c)$ the \textbf{superlevel set} \de{(Superniveaumenge,
obere Konturmenge)}. Written through a point $P$ rather than a level,
$W(P)=W(f(P))$ and $B(P)=B(f(P))$: everything at most as high as $P$, and
everything at least as high. The same sets appear in the optimization
volume as $L_{c}$ and $U_{c}$ --- lower and upper contour set --- and in
consumer theory $B(P)$ is the \emph{better set} of the bundle $P$.

The link to curvature, drilled later and worth seeing now:
\begin{itemize}[leftmargin=1.6em,itemsep=0.3ex]
  \item $f$ convex $\Longrightarrow$ every $W(c)$ is convex;
  \item $f$ concave $\Longrightarrow$ every $B(c)$ is convex;
  \item every $B(c)$ convex is the \emph{definition} of
        \textbf{quasiconcave}, which is weaker: $f(x,y)=xy$ on the positive
        quadrant is quasiconcave but not concave.
\end{itemize}
\end{rulebox}

\begin{gallery}
\begin{tikzpicture}
  \begin{axis}[geoaxis, xmin=-3.2, xmax=3.2, ymin=-3.2, ymax=3.2,
      xtick={-2,2}, ytick={-2,2}, title={$W(4)$ for $f=x^{2}+y^{2}$}]
    \addplot[draw=none, name path=up, domain=-2:2, samples=201]{sqrt(4-x^2)};
    \addplot[draw=none, name path=lo, domain=-2:2, samples=201]{-sqrt(4-x^2)};
    \region{up}{lo}{}
    \addplot[domain=0:360, samples=181]({2*cos(x)},{2*sin(x)});
  \end{axis}
\end{tikzpicture}\hfill
\begin{tikzpicture}
  \begin{axis}[geoaxis, xmin=-0.4, xmax=8.4, ymin=-0.4, ymax=8.4,
      xtick={2,4,6,8}, ytick={2,4,6,8}, title={$B(4)$ for $f=xy$, $x,y>0$}]
    \addplot[draw=none, name path=hyp, domain=0.5:8, samples=201]{4/x};
    \addplot[draw=none, name path=top, domain=0.5:8, samples=2]{8.4};
    \region{hyp}{top}{}
    \addplot[domain=0.5:8, samples=201]{4/x};
  \end{axis}
\end{tikzpicture}
\gcap{A sublevel set of a convex function and a superlevel set of a
quasiconcave one. Both are convex sets; the second is the \emph{better set}
of a Cobb--Douglas bundle.}
\end{gallery}

\begin{rulebox}{The standard shapes}
Six surfaces cover most of what appears in exercises, and each has a
signature contour map.
\begin{itemize}[leftmargin=1.6em,itemsep=0.3ex]
  \item \textbf{Plane} $f=ax+by+c$: contours are \emph{parallel straight
        lines}, equally spaced. A budget line is one of them.
  \item \textbf{Paraboloid} $f=x^{2}+y^{2}$: a bowl; contours are
        concentric circles of radius $\sqrt{c}$, so they get
        \emph{closer together} outwards --- the bowl steepens.
  \item \textbf{Cone} $f=\sqrt{x^{2}+y^{2}}$: contours are concentric
        circles of radius $c$, \emph{equally spaced} --- constant slope.
  \item \textbf{Saddle} $f=x^{2}-y^{2}$: contours are hyperbolas; the level
        $c=0$ degenerates into the crossing pair $y=\pm x$.
  \item \textbf{Rotated saddle} $f=xy$: hyperbolas again, but with the
        \emph{axes} as asymptotes. Restricted to the positive quadrant this
        is the Cobb--Douglas picture.
  \item \textbf{Parabolic cylinder} $f=x^{2}$: the surface is a trough,
        constant in $y$; contours are \emph{pairs of vertical lines}
        $x=\pm\sqrt{c}$.
\end{itemize}
Circles alone do not identify the surface --- bowl, cone and dome all have
circular contours. The \emph{spacing} and the direction in which the labels
grow settle it.
\end{rulebox}

\begin{rulebox}{Reading curvature off the picture}
\begin{itemize}[leftmargin=1.6em,itemsep=0.3ex]
  \item \textbf{Convex} (bowl-shaped): every chord drawn between two points
        of the surface lies \emph{above} it; no part of the surface bulges
        over a chord. All sublevel sets are convex.
  \item \textbf{Concave} (dome-shaped): every chord lies \emph{below} the
        surface. All superlevel sets are convex.
  \item \textbf{Neither}: a saddle. Some chords lie above, some below.
  \item \textbf{Quasiconcave without being concave}: the superlevel sets
        are still convex --- nested ``better sets'' bulging towards the
        origin --- while the surface itself is not a dome. The Cobb--Douglas
        hill is the standard example.
  \item A \textbf{linear} function is both convex and concave, and both
        quasiconvex and quasiconcave.
\end{itemize}
\end{rulebox}

\begin{rulebox}{Economic readings}
\begin{itemize}[leftmargin=1.6em,itemsep=0.3ex]
  \item An \textbf{isoquant} is a level curve of a production function:
        all input combinations yielding the same output. An
        \textbf{indifference curve} is a level curve of utility.
  \item The slope of a level curve is the \textbf{marginal rate of
        substitution}, computed implicitly as $-f_{x}/f_{y}$
        (Chapter~5). Convex-to-the-origin level curves mean a diminishing
        rate of substitution.
  \item A \textbf{budget set} $\{p_{x}x+p_{y}y\le M,\ x,y\ge0\}$ is the
        sublevel set of a linear function intersected with the quadrant ---
        a triangle, and convex.
  \item The consumer's problem is to reach the highest indifference curve
        that still touches the budget triangle: the contour map plus the
        shaded set is the whole picture.
\end{itemize}
\end{rulebox}

\section{Worked Example}

\begin{worked}{Read the topography of $f(x,y)=x^{2}+4y^{2}$ and of $g(x,y)=x^{2}-y^{2}$: level curves, sub- and superlevel sets, and the surfaces}
\wstep{Step 1. A level curve is what you get by \emph{fixing the height}.}
Set $f(x,y)=c$ and solve for the curve in the plane:
\[
  x^{2}+4y^{2}=c .
\]
For $c>0$ divide by $c$ to reach the standard form of an ellipse:
\[
  \frac{x^{2}}{c}+\frac{y^{2}}{c/4}=1,
  \qquad\text{semi-axes } \sqrt c \text{ and } \tfrac{\sqrt c}{2} .
\]
So the contour map is a family of concentric ellipses, twice as wide as they
are tall. For $c=0$ the level ``curve'' collapses to the single point
$(0,0)$; for $c<0$ it is empty, since $f\ge0$ always.
\wtrap{Check which values of $c$ occur before drawing. Level sets can be
empty, a single point, or several disconnected pieces --- $x^{2}+4y^{2}=-1$
has no solutions, and a level set of $x^{2}-y^{2}$ has two branches.}

\wstep{Step 2. Spacing carries the steepness.}
The contours for $c=1,2,3,4$ have semi-axes $1,\sqrt2,\sqrt3,2$ in $x$. The
gaps shrink as you move out --- $1,\ 0.41,\ 0.32,\ 0.27$ --- so the surface
gets \emph{steeper} away from the origin. Contours that crowd together mean
a steep slope; contours far apart mean a plateau. That is the one skill a
contour map is for.

\wstep{Step 3. Sub- and superlevel sets.}
\[
  \{f\le c\}=\text{the filled ellipse: convex},
  \qquad
  \{f\ge c\}=\text{everything outside it: not convex.}
\]
$f$ is a convex function, so its \textbf{sub}level sets are convex. For a
concave function such as a utility $U$, it is the \textbf{super}level sets
$\{U\ge\bar u\}$ --- the ``at least as good as'' sets --- that are convex,
and that is precisely the convexity of preferences.
\wtrap{Convex function and convex set are different words. Match them
carefully: convex function $\Rightarrow$ sublevel sets convex; concave
function $\Rightarrow$ superlevel sets convex. The direction is chosen by
the inequality, not by the word ``convex''.}

\wstep{Step 4. The surface itself.}
$z=x^{2}+4y^{2}$ is an \textbf{elliptic paraboloid}: a bowl, with its
minimum at the origin, rising in every direction, steeper along $y$ than
along $x$ (the coefficient $4$ against $1$).

\wstep{Step 5. Now $g(x,y)=x^{2}-y^{2}$, where the sign flips.}
\[
  g=c:\quad x^{2}-y^{2}=c ,
\]
a family of \textbf{hyperbolas} --- opening left--right for $c>0$, up--down
for $c<0$. The level set $c=0$ degenerates into the \emph{pair of lines}
$y=\pm x$, which is where the two families change over.

\wstep{Step 6. Two pictures, side by side.}
\begin{center}
\begin{tikzpicture}
\begin{axis}[contouraxis, width=5.2cm, height=5.2cm, xmin=-2.6, xmax=2.6,
             ymin=-2.6, ymax=2.6, title={$x^{2}+4y^{2}=c$}]
  \foreach \c in {0.5,1,2,3,4} {
    \addplot[domain=0:360, samples=161, variable=\t]
      ({sqrt(\c)*cos(\t)}, {sqrt(\c)/2*sin(\t)});
  }
\end{axis}
\end{tikzpicture}
\hfill
\begin{tikzpicture}
\begin{axis}[contouraxis, width=5.2cm, height=5.2cm, xmin=-2.6, xmax=2.6,
             ymin=-2.6, ymax=2.6, title={$x^{2}-y^{2}=c$}]
  \foreach \c in {0.5,2} {
    \addplot[domain=-1.6:1.6, samples=121] ({sqrt(\c)*cosh(x)}, {sqrt(\c)*sinh(x)});
    \addplot[domain=-1.6:1.6, samples=121] ({-sqrt(\c)*cosh(x)}, {sqrt(\c)*sinh(x)});
    \addplot[domain=-1.6:1.6, samples=121] ({sqrt(\c)*sinh(x)}, {sqrt(\c)*cosh(x)});
    \addplot[domain=-1.6:1.6, samples=121] ({sqrt(\c)*sinh(x)}, {-sqrt(\c)*cosh(x)});
  }
  \addplot[gray!60, thin, samples=2, domain=-2.6:2.6, forget plot] {x};
  \addplot[gray!60, thin, samples=2, domain=-2.6:2.6, forget plot] {-x};
\end{axis}
\end{tikzpicture}
\end{center}
\begin{center}
\begin{tikzpicture}
\begin{axis}[surfaxis, title={$z=x^{2}+4y^{2}$}]
  \addplot3[surf, domain=-2:2, y domain=-1:1] {x^2+4*y^2};
\end{axis}
\end{tikzpicture}
\hfill
\begin{tikzpicture}
\begin{axis}[surfaxis, title={$z=x^{2}-y^{2}$}]
  \addplot3[surf, domain=-2:2, y domain=-2:2] {x^2-y^2};
\end{axis}
\end{tikzpicture}
\end{center}

\wstep{Step 7. Read the origin off each map.}
For $f$ the contours \emph{close around} the origin, so it is a minimum: go
any direction and you climb. For $g$ the contours \emph{cross} at the
origin, and the crossing is the signature of a \textbf{saddle}: along the
$x$-axis $g=x^{2}$ rises, along the $y$-axis $g=-y^{2}$ falls. The point is
a minimum in one direction and a maximum in another, so it is neither.
\wtrap{Closed loops shrinking to a point mean an extremum; crossing curves
mean a saddle. Checking only one direction --- for instance only along the
axes of $g$ where the answers disagree --- is exactly how a saddle gets
misreported as an optimum in Volume~7.}

\wstep{Step 8. The economic vocabulary.}
Level curves of a utility function are \textbf{indifference curves}; of a
production function, \textbf{isoquants}; of a cost function,
\textbf{isocost lines}. Their slope is the marginal rate of substitution,
computed in Chapter~5 by implicit differentiation, and the standard
``convex to the origin'' shape is exactly the convexity of the superlevel
sets in Step~3.
\end{worked}

\section{Practice Problems}

State the domain of each function as a region in the plane, in problems
1--8, and describe it in words.

\begin{problems}
  \item $f(x,y)=\sqrt{16-x^{2}-y^{2}}$
  \item $f(x,y)=\ln(x-y)$
  \item $f(x,y)=\dfrac{1}{x^{2}+y^{2}-1}$
  \item $f(x,y)=\dfrac{xy}{x-y}$
  \item $f(x,y)=\sqrt{x}+\sqrt{y}$
  \item $f(x,y)=\ln\bigl(xy\bigr)$
  \item $f(x,y)=\sqrt{x^{2}+y^{2}-4}$
  \item $f(K,L)=K^{0.4}L^{0.6}$, and state the economic domain too.
\end{problems}

In problems 9--20, give the level set $N_{c}$ for the stated level, solved
for $y$ where possible, and name the curve.

\begin{problems}[resume]
  \item $f(x,y)=x+2y$, \;$c=4$.
  \item $f(x,y)=x^{2}+y^{2}$, \;$c=9$.
  \item $f(x,y)=x^{2}+y^{2}$, \;$c=0$.
  \item $f(x,y)=x^{2}+4y^{2}$, \;$c=4$.
  \item $f(x,y)=xy$, \;$c=6$.
  \item $f(x,y)=x^{2}-y^{2}$, \;$c=0$.
  \item $f(x,y)=x^{2}-y^{2}$, \;$c=4$.
  \item $f(x,y)=y-x^{2}$, \;$c=2$.
  \item $f(x,y)=\sqrt{x^{2}+y^{2}}$, \;$c=3$.
  \item $f(x,y)=\mathrm{e}^{x+y}$, \;$c=1$.
  \item $f(K,L)=\sqrt{KL}$, \;$c=10$.
  \item $f(x,y)=2x+3y$, \;$c=12$; and say what this curve is called when
        $f$ is a cost.
\end{problems}

Problems 21--28 refer to the contour map below, which shows the level
curves of $f(x,y)=x^{3}-3x+y^{2}$ at the levels $c=-1,0,2,4$, together with
six marked points.

\begin{gallery}
\hspace*{\fill}%
\begin{tikzpicture}
  \begin{axis}[contouraxis, width=8.2cm, height=8.2cm,
      xmin=-2.6, xmax=2.6, ymin=-3.2, ymax=3.2,
      xtick={-2,-1,1,2}, ytick={-2,-1,1,2}]
    % c = 4
    \addplot[domain=-2.2:2.196, samples=201]{ sqrt(4-x^3+3*x)};
    \addplot[domain=-2.2:2.196, samples=201]{-sqrt(4-x^3+3*x)};
    % c = 2 (through the saddle at (-1,0))
    \addplot[domain=-2.2:2, samples=201]{ sqrt(2-x^3+3*x)};
    \addplot[domain=-2.2:2, samples=201]{-sqrt(2-x^3+3*x)};
    % c = 0, two pieces
    \addplot[domain=-2.2:-1.7321, samples=101]{ sqrt(-x^3+3*x)};
    \addplot[domain=-2.2:-1.7321, samples=101]{-sqrt(-x^3+3*x)};
    \addplot[domain=0:1.7321, samples=101]{ sqrt(-x^3+3*x)};
    \addplot[domain=0:1.7321, samples=101]{-sqrt(-x^3+3*x)};
    % c = -1, two pieces
    \addplot[domain=-2.2:-1.8794, samples=61]{ sqrt(-1-x^3+3*x)};
    \addplot[domain=-2.2:-1.8794, samples=61]{-sqrt(-1-x^3+3*x)};
    \addplot[domain=0.3473:1.5321, samples=101]{ sqrt(-1-x^3+3*x)};
    \addplot[domain=0.3473:1.5321, samples=101]{-sqrt(-1-x^3+3*x)};
    % marked points
    \addplot[only marks, mark=*, mark size=1.4pt, black, forget plot]
      coordinates {(0,0) (1,0) (-1,0) (2,1) (0,2) (-2,1)};
    \node[font=\tiny, anchor=north east] at (axis cs:-0.06,-0.12) {$P_{1}$};
    \node[font=\tiny, anchor=north west] at (axis cs:1.08,-0.12) {$P_{2}$};
    \node[font=\tiny, anchor=south east] at (axis cs:-1.12,0.12) {$P_{3}$};
    \node[font=\tiny, anchor=south west] at (axis cs:2.05,1.05) {$P_{4}$};
    \node[font=\tiny, anchor=south west] at (axis cs:0.05,2.05) {$P_{5}$};
    \node[font=\tiny, anchor=south east] at (axis cs:-2.05,1.05) {$P_{6}$};
  \end{axis}
\end{tikzpicture}%
\hspace*{\fill}
\gcap{Level curves of $f(x,y)=x^{3}-3x+y^{2}$ at $c=-1,0,2,4$.}
\end{gallery}

\begin{problems}[resume]
  \item Give $f(P_{i})$ for all six marked points
        $P_{1}=(0,0)$, $P_{2}=(1,0)$, $P_{3}=(-1,0)$, $P_{4}=(2,1)$,
        $P_{5}=(0,2)$, $P_{6}=(-2,1)$.
  \item Which of the six points lie on one and the same level curve?
  \item One of the drawn contours crosses itself. At which point, and what
        does that tell you about the surface there?
  \item One of the drawn contours has shrunk to a single point. Which
        level, which point, and what feature of the surface is it?
  \item From the picture alone, name the two critical points of $f$ and
        classify each as a local maximum, a local minimum, or a saddle.
  \item Confirm the previous answer by computation: solve
        $f_{x}=f_{y}=0$.
  \item Is the sublevel set $W(-1)$ connected? Answer from the picture and
        explain what the two pieces mean.
  \item Moving away from $P_{2}$, in which direction does $f$ increase
        fastest? Answer qualitatively from the contours, then check with
        $\grad f$ at a nearby point, say $(1.5,0.5)$.
\end{problems}

\begin{problems}[resume]
  \item Match each surface (a)--(f) below to its contour map (1)--(6).
  \item For each of the six functions in the previous problem, say whether
        it is convex, concave, both, or neither on all of $\R^{2}$.
\end{problems}

\begin{gallery}
\begin{tikzpicture}
  \begin{axis}[surfaxis, title={(a)}]
    \addplot3[surf, domain=-2:2, y domain=-2:2]{x+y};
  \end{axis}
\end{tikzpicture}\hfill
\begin{tikzpicture}
  \begin{axis}[surfaxis, title={(b)}, z post scale=1.2]
    \addplot3[surf, domain=-2:2, y domain=-2:2]{x^2+y^2};
  \end{axis}
\end{tikzpicture}\hfill
\begin{tikzpicture}
  \begin{axis}[surfaxis, title={(c)}, z post scale=1.2]
    \addplot3[surf, domain=-2:2, y domain=-2:2]{sqrt(x^2+y^2)};
  \end{axis}
\end{tikzpicture}
\end{gallery}

\begin{gallery}
\begin{tikzpicture}
  \begin{axis}[surfaxis, title={(d)}, z post scale=1.2]
    \addplot3[surf, domain=-2:2, y domain=-2:2]{x^2-y^2};
  \end{axis}
\end{tikzpicture}\hfill
\begin{tikzpicture}
  \begin{axis}[surfaxis, title={(e)}, z post scale=1.2, view={105}{32}]
    \addplot3[surf, domain=-2:2, y domain=-2:2]{x*y};
  \end{axis}
\end{tikzpicture}\hfill
\begin{tikzpicture}
  \begin{axis}[surfaxis, title={(f)}, z post scale=1.2]
    \addplot3[surf, domain=-2:2, y domain=-2:2]{x^2};
  \end{axis}
\end{tikzpicture}
\end{gallery}

\begin{gallery}
\begin{tikzpicture}
  \begin{axis}[contouraxis, xmin=-2.4, xmax=2.4, ymin=-2.4, ymax=2.4,
      xtick={-2,2}, ytick={-2,2}, title={(1)}]
    \foreach \c in {-2,-1,0,1,2} { \addplot[domain=-2.3:2.3]{\c-x}; }
  \end{axis}
\end{tikzpicture}\hfill
\begin{tikzpicture}
  \begin{axis}[contouraxis, xmin=-2.4, xmax=2.4, ymin=-2.4, ymax=2.4,
      xtick={-2,2}, ytick={-2,2}, title={(2)}]
    \foreach \r in {0.25,0.5,0.75,1,1.25,1.5,1.75,2} {
      \addplot[domain=0:360, samples=121]({\r*cos(x)},{\r*sin(x)}); }
  \end{axis}
\end{tikzpicture}\hfill
\begin{tikzpicture}
  \begin{axis}[contouraxis, xmin=-2.4, xmax=2.4, ymin=-2.4, ymax=2.4,
      xtick={-2,2}, ytick={-2,2}, title={(3)}]
    \foreach \r in {0.707,1,1.225,1.414,1.581,1.732,1.871,2} {
      \addplot[domain=0:360, samples=121]({\r*cos(x)},{\r*sin(x)}); }
  \end{axis}
\end{tikzpicture}
\end{gallery}

\begin{gallery}
\begin{tikzpicture}
  \begin{axis}[contouraxis, xmin=-2.4, xmax=2.4, ymin=-2.4, ymax=2.4,
      xtick={-2,2}, ytick={-2,2}, title={(4)}]
    \addplot[domain=-2.3:2.3]{x};
    \addplot[domain=-2.3:2.3]{-x};
    \foreach \s in {1,-1} {
      \addplot[domain=-1.4:1.4, samples=81]({\s*cosh(x)},{sinh(x)});
      \addplot[domain=-1.4:1.4, samples=81]({sinh(x)},{\s*cosh(x)}); }
  \end{axis}
\end{tikzpicture}\hfill
\begin{tikzpicture}
  \begin{axis}[contouraxis, xmin=-2.4, xmax=2.4, ymin=-2.4, ymax=2.4,
      xtick={-2,2}, ytick={-2,2}, title={(5)}]
    \foreach \c in {1,2} {
      \addplot[domain=0.42:2.3, samples=101]{\c/x};
      \addplot[domain=-2.3:-0.42, samples=101]{\c/x};
      \addplot[domain=0.42:2.3, samples=101]{-\c/x};
      \addplot[domain=-2.3:-0.42, samples=101]{-\c/x}; }
  \end{axis}
\end{tikzpicture}\hfill
\begin{tikzpicture}
  \begin{axis}[contouraxis, xmin=-2.4, xmax=2.4, ymin=-2.4, ymax=2.4,
      xtick={-2,2}, ytick={-2,2}, title={(6)}]
    \foreach \v in {-2,-1.414,-1,1,1.414,2} {
      \addplot[samples=2] coordinates {(\v,-2.3) (\v,2.3)}; }
  \end{axis}
\end{tikzpicture}
\end{gallery}

In problems 31--38, describe the stated sub- or superlevel set as a region,
and say whether it is convex.

\begin{problems}[resume]
  \item $W(4)$ for $f(x,y)=x^{2}+y^{2}$.
  \item $B(4)$ for $f(x,y)=x^{2}+y^{2}$.
  \item $W(6)$ for $f(x,y)=2x+3y$.
  \item $B(4)$ for $f(x,y)=xy$ restricted to $x,y>0$.
  \item $W(0)$ for $f(x,y)=x^{2}-y^{2}$.
  \item $B(1)$ for $f(x,y)=\sqrt{x^{2}+y^{2}}$.
  \item $W(P)$ for $f(x,y)=x^{2}+y^{2}$ with $P=(3,4)$.
  \item $B(P)$ for $f(K,L)=\sqrt{KL}$ with $P=(4,9)$, on $K,L>0$.
\end{problems}

\begin{problems}[resume]
  \item Sketch the level sets $W(2)$, $B(2)$, $W(-1)$ and $B(-1)$ of the
        one-variable function $g(x)=x^{3}-3x$ as subsets of the $x$-axis,
        and say which of the four are convex sets.
  \item Explain why for a function of one variable a level set
        $W(c)$ is convex exactly when it is an interval.
  \item A firm's isoquants are the curves $KL=\bar Q$. Sketch three of them
        for $\bar Q=1,4,9$ on $K,L>0$, and say whether the isoquants get
        closer together or further apart as output rises along the ray
        $K=L$.
  \item For the budget set $\{2x+5y\le100,\ x,y\ge0\}$: give the two
        intercepts, sketch the set, and say why it is convex.
  \item The indifference curves of $U(x,y)=xy$ and the budget line
        $2x+5y=100$ are drawn on the same axes. Describe, in words and
        without computing, the geometric condition that locates the optimal
        bundle.
  \item Two contour maps show circles as level curves. In the first the
        circles are equally spaced; in the second they crowd together
        outwards. Which is the cone and which the paraboloid, and why?
  \item A contour map has no closed loops and no crossings anywhere, only
        parallel straight lines. What can be said about the function?
\end{problems}

\section{Solutions}

\begin{solutions}
  \item $16-x^{2}-y^{2}\ge0\iff x^{2}+y^{2}\le16$: the closed disc of
        radius $4$ about the origin.
  \item $x-y>0$: the open half-plane strictly below the line $y=x$.
  \item $x^{2}+y^{2}\neq1$: the whole plane except the unit circle.
  \item $x\neq y$: the plane without the line $y=x$.
  \item $x\ge0$ and $y\ge0$: the closed non-negative quadrant.
  \item $xy>0$: the open first and third quadrants (the axes excluded).
  \item $x^{2}+y^{2}\ge4$: everything outside or on the circle of radius
        $2$ --- the complement of an open disc.
  \item The exponents are not integers, so $K\ge0$ and $L\ge0$: the closed
        non-negative quadrant. Economically the same set, usually taken
        open ($K,L>0$) so that marginal products exist.
  \item $y=\dfrac{4-x}{2}=2-\tfrac12x$: a straight line of slope
        $-\tfrac12$. All level sets are parallel lines.
  \item $x^{2}+y^{2}=9$, i.e. $y=\pm\sqrt{9-x^{2}}$: the circle of radius
        $3$.
  \item $x^{2}+y^{2}=0$ forces $x=y=0$: the level set is the single point
        $(0,0)$, the bottom of the bowl.
  \item $\dfrac{x^{2}}{4}+y^{2}=1$, so $y=\pm\sqrt{1-x^{2}/4}$: an ellipse
        with semi-axes $2$ and $1$.
  \item $y=6/x$: a hyperbola in the first and third quadrants, with both
        axes as asymptotes.
  \item $x^{2}=y^{2}\iff y=\pm x$: the \emph{two crossing lines} through
        the origin. This degenerate level is the signature of a saddle.
  \item $x^{2}-y^{2}=4$, so $y=\pm\sqrt{x^{2}-4}$, defined for
        $|x|\ge2$: a hyperbola opening left and right.
  \item $y=x^{2}+2$: a parabola, the standard parabola shifted up by $2$.
        All level sets are vertical translates of it.
  \item $\sqrt{x^{2}+y^{2}}=3\iff x^{2}+y^{2}=9$: the circle of radius
        $3$. Note that for the cone the radius equals $c$, not
        $\sqrt{c}$.
  \item $\mathrm{e}^{x+y}=1\iff x+y=0$, so $y=-x$: a straight line. Every
        level set of $\mathrm{e}^{x+y}$ is a line $x+y=\ln c$, defined only
        for $c>0$.
  \item $\sqrt{KL}=10\iff KL=100\iff L=100/K$: a hyperbola --- the
        isoquant for output $10$.
  \item $y=\dfrac{12-2x}{3}=4-\tfrac23x$: a straight line. For a cost
        function this is an \textbf{isocost line}; the analogous level
        curve of a budget constraint is the budget line.
  \item $f(x,y)=x^{3}-3x+y^{2}$ gives
        $f(P_{1})=0$, \;$f(P_{2})=1-3+0=-2$, \;$f(P_{3})=-1+3+0=2$,
        \;$f(P_{4})=8-6+1=3$, \;$f(P_{5})=0-0+4=4$,
        \;$f(P_{6})=-8+6+1=-1$.
  \item None of them share a level: the six values
        $0,-2,2,3,4,-1$ are distinct. $P_{3}$ lies on the drawn contour
        $c=2$, $P_{5}$ on $c=4$ and $P_{6}$ on $c=-1$; $P_{1}$, $P_{2}$ and
        $P_{4}$ lie between drawn contours.
  \item The contour $c=2$ crosses itself at $P_{3}=(-1,0)$. A
        self-crossing level curve marks a \textbf{saddle point}: from
        $(-1,0)$ the surface rises in two opposite directions and falls in
        the two others.
  \item The level $c=-2$ is a single point, $(1,0)$. Nested contours
        shrinking to a point mark a \textbf{local extremum}; since the
        levels decrease inwards, it is a local \emph{minimum}, with value
        $-2$.
  \item From the picture: a local minimum at $(1,0)$ (contours closing in
        on it) and a saddle at $(-1,0)$ (the crossing contour). There is no
        local maximum --- no contours close in on a point with the labels
        increasing inwards.
  \item $f_{x}=3x^{2}-3=0\Rightarrow x=\pm1$; $f_{y}=2y=0\Rightarrow y=0$.
        Critical points $(1,0)$ and $(-1,0)$. With $f_{xx}=6x$,
        $f_{yy}=2$, $f_{xy}=0$: at $(1,0)$ the determinant is $12>0$ with
        $f_{xx}>0$, a local minimum; at $(-1,0)$ it is $-12<0$, a saddle.
  \item No: $W(-1)$ has two pieces --- a small oval around the minimum
        $(1,0)$ and an unbounded piece on the far left. A sublevel set of a
        non-convex function may fall apart into components, which is
        exactly what convexity of $f$ would forbid.
  \item Away from $P_{2}=(1,0)$ the contours are closed ovals around it, so
        $f$ increases in \emph{every} direction, fastest where the ovals
        are most tightly packed --- to the right. Checking,
        $\grad f=(3x^{2}-3,\,2y)$, so at $(1.5,0.5)$,
        $\grad f=(3.75,\,1)$: pointing right and slightly up, and
        perpendicular to the contour there.
  \item (a)$\to$(1) plane, parallel lines. \;(b)$\to$(3) paraboloid:
        circles of radius $\sqrt{c}$, crowding outwards.
        \;(c)$\to$(2) cone: circles of radius $c$, equally spaced.
        \;(d)$\to$(4) saddle $x^{2}-y^{2}$: hyperbolas with the diagonals
        $y=\pm x$ as the degenerate level. \;(e)$\to$(5) $xy$: hyperbolas
        with the \emph{axes} as asymptotes. \;(f)$\to$(6) parabolic
        cylinder: pairs of vertical lines.
  \item (a) linear, hence \textbf{both} convex and concave. (b) convex (a
        bowl). (c) convex (a cone opening upwards; not differentiable at
        the origin, but still convex). (d) neither --- a saddle.
        (e) neither on $\R^{2}$; on the positive quadrant it is
        quasiconcave. (f) convex: $x^{2}$ is convex and adding no
        $y$-dependence preserves it.
  \item The closed disc $x^{2}+y^{2}\le4$ of radius $2$. Convex.
  \item The complement of the open disc, $x^{2}+y^{2}\ge4$: everything
        outside the circle. \textbf{Not} convex --- the segment joining
        $(-2,0)$ to $(2,0)$ passes through the origin, which is not in the
        set.
  \item The half-plane $2x+3y\le6$, bounded by the line through $(3,0)$ and
        $(0,2)$. Convex: every half-plane is.
  \item $\{(x,y): x,y>0,\ xy\ge4\}$, the region on the far side of the
        hyperbola $y=4/x$. Convex --- this is precisely the statement that
        $xy$ is quasiconcave on the quadrant.
  \item $x^{2}-y^{2}\le0\iff|y|\ge|x|$: the two opposite ``vertical'' sectors
        between the lines $y=\pm x$. Not convex: $(0,1)$ and $(0,-1)$ lie
        in it, but the midpoint $(0,0)$ is on the boundary and points just
        off the segment leave the set --- more simply, the set is not even
        connected in the way a convex set must be, since joining
        $(0,1)$ to $(1,2)$ stays inside but joining $(0,1)$ to $(0,-1)$
        passes through the pinch at the origin only.
  \item $\sqrt{x^{2}+y^{2}}\ge1$: everything outside or on the unit circle.
        Not convex, for the same reason as problem 32.
  \item $f(P)=9+16=25$, so $W(P)=\{x^{2}+y^{2}\le25\}$, the closed disc of
        radius $5$ through $P$. Convex.
  \item $f(P)=\sqrt{36}=6$, so
        $B(P)=\{(K,L):K,L>0,\ \sqrt{KL}\ge6\}=\{KL\ge36\}$: the region
        beyond the isoquant $L=36/K$. Convex.
  \item $g(x)=x^{3}-3x$ has $g(2)=2$, $g(-1)=2$, $g(1)=-2$, $g(-2)=-2$.
        $W(2)=\{x\le2\}$ minus nothing --- solving $x^{3}-3x\le2$ gives
        $(x+1)^{2}(x-2)\le0$, i.e. $x\le2$: the ray $(-\infty,2]$,
        \textbf{convex}. $B(2)=[2,\infty)\cup\{-1\}$: the ray from $2$
        \emph{plus the isolated point} $x=-1$, \textbf{not convex}.
        $W(-1)$: solving $x^{3}-3x\le-1$ gives
        $(-\infty,-1.879]\cup[0.347,1.532]$, two intervals, \textbf{not
        convex}. $B(-1)=[-1.879,0.347]\cup[1.532,\infty)$, also two
        pieces, \textbf{not convex}.
  \item In $\R^{1}$ a set is convex exactly when it contains the whole
        segment between any two of its points, and a subset of the line
        with that property is precisely an interval (possibly unbounded or
        empty). So ``convex'' and ``interval'' are the same statement one
        dimension down --- which is why the one-variable picture is the
        right place to first meet the definition.
  \item $L=\bar Q/K$: three hyperbola branches, $L=1/K$, $L=4/K$ and
        $L=9/K$. Along the ray $K=L$ they cut it at
        $K=1,2,3$: the isoquants for equally spaced \emph{square roots} of
        output, so for equally spaced output levels they move
        \emph{further apart}. Rising output requires ever larger input
        increases --- diminishing returns along the ray.
  \item Intercepts $(50,0)$ and $(0,20)$. The set is the triangle with
        vertices $(0,0)$, $(50,0)$, $(0,20)$. It is convex because it is an
        intersection of three half-planes ($2x+5y\le100$, $x\ge0$,
        $y\ge0$), and any intersection of convex sets is convex.
  \item The optimum is where the budget line is \textbf{tangent} to an
        indifference curve: the highest curve of the family that still
        touches the triangle touches it at exactly one point, and there the
        two curves have the same slope. Crossing would mean a higher curve
        is still reachable; missing the set entirely means the bundle is
        unaffordable.
  \item Equally spaced circles are the \textbf{cone}
        $\sqrt{x^{2}+y^{2}}$, whose level $c$ is a circle of radius $c$ ---
        constant slope, so equal height steps need equal horizontal steps.
        Circles crowding outwards are the \textbf{paraboloid}
        $x^{2}+y^{2}$, level $c$ having radius $\sqrt{c}$: the surface gets
        steeper, so equal height steps need ever smaller horizontal steps.
  \item Parallel straight contours, equally spaced or not, mean the
        function depends on the two variables only through one linear
        combination: $f(x,y)=h(ax+by)$ for some one-variable $h$. If the
        spacing is also constant, $h$ is linear and the surface is a plane.
        There are no critical points of extremum type anywhere.
\end{solutions}

\section{Additional Problems}

\begin{problems}
  \item Domain of $f(x,y)=\dfrac{\sqrt{y-x^{2}}}{x}$, as a region.
  \item Domain of $f(x,y)=\ln(4-x^{2}-y^{2})$, as a region.
  \item Domain and range of $f(x,y)=\mathrm{e}^{-(x^{2}+y^{2})}$.
  \item Level set of $f(x,y)=x^{2}+y^{2}-4x$ at $c=0$; name and locate the
        curve.
  \item Level set of $f(x,y)=x^{2}+y^{2}-4x$ at $c=-4$.
  \item Level set of $f(x,y)=y\,\mathrm{e}^{x}$ at $c=1$.
  \item Level set of $f(x,y)=\dfrac{y}{x}$ at $c=2$, and why the origin is
        excluded.
  \item Level sets of $f(x,y)=\max(x,y)$ at $c=1$: sketch and describe.
  \item For $f(x,y)=x^{2}+4y^{2}$, sketch the contours $c=4,16,36$ and say
        whether they crowd inwards or outwards.
  \item For $f(x,y)=x-y^{2}$, give the contours at $c=-1,0,1$ and describe
        the family.
  \item Give a function whose level curves are concentric circles that
        crowd \emph{inwards} towards the origin.
  \item Give a function of two variables whose level sets are all straight
        lines through the origin.
  \item $W(1)$ and $B(1)$ for $f(x,y)=|x|+|y|$: describe both and say which
        is convex.
  \item $W(0)$ for $f(x,y)=xy$ on all of $\R^{2}$: describe it and decide
        convexity.
  \item $B(2)$ for $f(x,y)=x+y$: describe it and decide convexity.
  \item $W(4)$ for $f(x,y)=x^{2}$: describe it and decide convexity.
  \item For $f(x,y)=x^{2}+y^{2}$ and $P=(1,2)$, give $W(P)$ and $B(P)$ and
        the level curve through $P$.
  \item A function has convex sublevel sets \emph{and} convex superlevel
        sets at every level. What kind of function can it be? Give an
        example.
  \item Sketch the surface $f(x,y)=4-x^{2}-y^{2}$ in words, give its
        contours, and classify it as convex or concave.
  \item Two functions have exactly the same family of level curves but
        different surfaces. Give such a pair and explain how the contour
        maps can nevertheless be told apart.
  \item Classify $f(x,y)=x^{2}+y^{2}$, $g(x,y)=-x^{2}-y^{2}$,
        $h(x,y)=x^{2}-y^{2}$ and $k(x,y)=3x-2y$ as convex, concave, both or
        neither.
  \item Which of the four functions in the previous problem are
        quasiconcave?
  \item The isoquants of a production function are straight lines with
        slope $-1$. What does that say about substitutability of the two
        inputs, and give such a production function.
  \item An indifference map consists of circles centred at $(4,4)$, with
        utility \emph{increasing} inwards. Interpret the point $(4,4)$
        economically and say why such preferences are unusual.
  \item On the contour map of problem 21, a walker starts at $P_{6}=(-2,1)$
        and always walks perpendicular to the contours, downhill. Describe
        where they end up and why the answer is not the minimum at
        $(1,0)$.
\end{problems}

\section{Additional Solutions}

\begin{solutions}
  \item $y-x^{2}\ge0$ and $x\neq0$: the region on or above the parabola
        $y=x^{2}$, with the $y$-axis removed.
  \item $4-x^{2}-y^{2}>0\iff x^{2}+y^{2}<4$: the \emph{open} disc of
        radius $2$ (boundary excluded, because the logarithm needs a
        strictly positive argument).
  \item $D=\R^{2}$. The exponent $-(x^{2}+y^{2})$ runs over $(-\infty,0]$,
        so the range is $(0,1]$ --- a bell-shaped hill with peak $1$ at the
        origin.
  \item $x^{2}-4x+y^{2}=0\iff(x-2)^{2}+y^{2}=4$: the circle of radius $2$
        centred at $(2,0)$, passing through the origin.
  \item $(x-2)^{2}+y^{2}=0$: the single point $(2,0)$ --- the minimum of
        the shifted bowl, with value $-4$.
  \item $y\,\mathrm{e}^{x}=1\iff y=\mathrm{e}^{-x}$: a decaying
        exponential curve.
  \item $y/x=2\iff y=2x$ with $x\neq0$: the line of slope $2$ through the
        origin, \emph{minus} the origin itself, since $f$ is undefined
        where $x=0$.
  \item $\max(x,y)=1$ is the union of the ray $\{x=1,\ y\le1\}$ and the ray
        $\{y=1,\ x\le1\}$: an L-shaped corner opening down-left, with its
        vertex at $(1,1)$.
  \item Ellipses $\dfrac{x^{2}}{c}+\dfrac{4y^{2}}{c}=1$ with semi-axes
        $\sqrt{c}$ and $\sqrt{c}/2$: for $c=4,16,36$ the semi-axes are
        $2,4,6$ and $1,2,3$. Equally spaced, in fact --- but for
        \emph{equally spaced} levels the contours would crowd outwards, as
        for any quadratic bowl.
  \item $x=y^{2}+c$: parabolas opening to the right, the whole family being
        horizontal translates of $x=y^{2}$ by $c$. For $c=-1,0,1$ the
        vertices sit at $(-1,0)$, $(0,0)$, $(1,0)$.
  \item Any function whose height grows ever more slowly with the radius,
        for instance $f(x,y)=\sqrt[4]{x^{2}+y^{2}}$ or
        $f(x,y)=\ln(x^{2}+y^{2})$: the level $c$ has radius growing faster
        than linearly in $c$, so equal height steps take ever larger
        radial steps and the contours thin out --- equivalently they crowd
        towards the origin.
  \item $f(x,y)=y/x$ (defined for $x\neq0$): the level $c$ is the punctured
        line $y=cx$. Any function of the ratio $y/x$ alone has this
        property; such functions are exactly the homogeneous functions of
        degree $0$.
  \item $W(1)=\{|x|+|y|\le1\}$ is the closed diamond (square rotated
        $45^{\circ}$) with vertices $(\pm1,0)$, $(0,\pm1)$ --- convex.
        $B(1)$ is everything outside it --- not convex.
  \item $xy\le0$ is the union of the second and fourth quadrants. Not
        convex: $(-1,1)$ and $(1,-1)$ both lie in it, but their midpoint
        $(0,0)$ does lie in it too --- take instead $(-1,2)$ and $(2,-1)$,
        whose midpoint $(0.5,0.5)$ has $xy>0$ and is outside.
  \item $x+y\ge2$: a closed half-plane, convex.
  \item $x^{2}\le4\iff-2\le x\le2$: the vertical strip between the lines
        $x=-2$ and $x=2$, convex (an intersection of two half-planes).
  \item $f(P)=1+4=5$: the level curve is the circle of radius $\sqrt5$,
        $W(P)$ the closed disc of radius $\sqrt5$, and $B(P)$ everything
        outside or on that circle.
  \item If every sublevel \emph{and} every superlevel set is convex, the
        function is both quasiconvex and quasiconcave. The level sets must
        then be ``flat'': any function of the form $h(ax+by)$ with $h$
        monotone qualifies, the simplest being a linear function
        $f(x,y)=ax+by$.
  \item A dome: the paraboloid turned upside down, with its peak of height
        $4$ above the origin and falling away in every direction. Contours
        are circles $x^{2}+y^{2}=4-c$ of radius $\sqrt{4-c}$, shrinking as
        $c$ rises to $4$. The function is \textbf{concave}.
  \item $f(x,y)=x^{2}+y^{2}$ and $g(x,y)=\sqrt{x^{2}+y^{2}}$ both have
        concentric circles as level sets. The maps differ in the
        \emph{labels}: for equally spaced levels $f$'s contours crowd
        outwards while $g$'s stay equally spaced. A contour map without its
        level labels does not determine the function.
  \item $f$ convex; $g$ concave; $h$ neither (a saddle); $k$ linear, hence
        both.
  \item $f$: quasiconvex, not quasiconcave (its superlevel sets are
        complements of discs). $g$: concave, hence quasiconcave. $h$:
        neither. $k$: linear, hence both quasiconvex and quasiconcave.
  \item Straight isoquants of slope $-1$ mean the marginal rate of
        technical substitution is constant at $1$: the inputs are
        \textbf{perfect substitutes}, one unit of either always replacing
        one unit of the other. The production function is
        $Q(K,L)=K+L$ (or any increasing function of $K+L$).
  \item $(4,4)$ is a \textbf{bliss point}: a most-preferred bundle, beyond
        which more of either good makes the consumer worse off. Standard
        preferences are monotone --- more is better --- so utility should
        increase without bound in the direction of both goods; circular
        indifference curves with utility rising inwards describe satiation
        instead.
  \item Walking perpendicular to the contours downhill from $P_{6}=(-2,1)$
        means following $-\grad f$. At $P_{6}$,
        $\grad f=(3\cdot4-3,\,2)=(9,2)$, so the walker moves left and
        slightly down --- into the region $x<-2$, where the contours of
        \emph{lower} level continue without bound. The function
        $x^{3}-3x+y^{2}$ tends to $-\infty$ as $x\to-\infty$, so the
        descent never ends: the walker never reaches the local minimum at
        $(1,0)$, because gradient descent finds a \emph{local} basin only
        if it starts in one.
\end{solutions}
